\documentclass[11pt]{article}
\usepackage{fontspec}
\usepackage[french]{babel}
\usepackage{amsmath}
\usepackage{amssymb}
\usepackage{amsthm}
\usepackage{mathrsfs}
\usepackage{tikz-cd}
\usepackage{enumitem}
\frenchsetup{AutoSpacePunctuation=false}
\newcommand{\printedpage}[1]{\marginpar{\footnotesize p.~#1}}
\newcommand{\msp}[1]{\marginpar{\footnotesize [#1]}}
\begin{document}

\printedpage{78}
% head: A. Grothendieck : Dérivateurs, Ch. XII, Édition des Universités de Montpellier II et Paris VII
\begin{enumerate}
\item[b)] J'ai envie, pour des applications ultérieures, de disposer d'une classe d'immersions ouvertes ayant les propriétés suivantes :
\begin{enumerate}
\item[1°)] Ce sont des $W$-cofibrations, et
\item[2°)] Si $U \overset{i}{\hookrightarrow} X$ est dans la classe, et si $\varepsilon \xrightarrow{j} I$ est l'inclusion ‘source-but’ de la catégorie discrète $\{0,1\}$ dans l'intervalle
\[
\Phi_n = 0 \longrightarrow 1 \longleftarrow 2 \longrightarrow \dots 2n-1 \longleftarrow 2n
\]
$(j(0)=0$, $j(1)=2n)$, $n\geq 2$, l'inclusion $\alpha$ de $(U\times I)\cup(X\times\varepsilon)$ dans $X\times I$ est dans la classe.
\end{enumerate}
Or cette dernière propriété de stabilité n'est \textit{pas} vérifiée par les immersions de Dwyer. (Prendre $i : \{0\} \longrightarrow \boldsymbol{\Delta}^1 = \{0,1\}$, l'inclusion $\alpha$ n'est pas immersion de Dwyer.)
\end{enumerate}

\msp{page 115}

\subsection*{2.5\quad Complément : critère d'asphéricité totale d'une catégorie $W$-propre sur une autre.}

Soit
\[
f : X \longrightarrow I
\]
une catégorie propre sur une catégorie d'indices $I$. Si $I$ a un objet final $i_0$, alors $X/i_0$ s'identifie à $X$, et par suite l'inclusion
\[
X_{i_0} \longrightarrow X
\]
de la fibre de $X$ en $i_0$ est $W$-asphérique \texttt{[plutôt }$W$\texttt{-coasphérique]}, et a fortiori est dans $W$. Donc $X$ est asphérique si et seulement si $X_{i_0}$ l'est. Nous voulons étendre ce résultat en un critère de \textit{totale} $W$-asphéricité, de façon précise prouver la

\textbf{Proposition 11 :} \textit{Soit $f : X \longrightarrow I$ $W$-propre, et soit $i_0$ un objet final de $I$. Si $X_{i_0}$ est totalement $W$-asphérique} (déf. 1, page 88), \textit{il en est de même de $X$.} (Sous réserve que $W$ satisfasse l'axiome de connexité W(5).)

La démonstration sera donnée plus bas, après quelques préliminaires ayant un intérêt indépendant. Tout d'abord, rappelons que si $X$ est $W$-propre sur $I$, alors il donne naissance à un foncteur
\[
\begin{aligned}
I \quad &\longrightarrow \quad \mathrm{Hot}_W\\
i \quad &\longmapsto \quad \mathrm{hot}_W(X_i)(\simeq \mathrm{hot}_W(X/i));
\end{aligned}
\]
en le composant avec le foncteur
\[
\pi_0 : \mathrm{Hot}_W \longrightarrow \mathrm{Ens}
\]
(sous réserve que l'axiome de connexité W(5) (p. 10) soit satisfait)
% foot: Version 7 décembre 2002
\printedpage{79}
% head: A. Grothendieck : Dérivateurs, Ch. XII, Édition des Universités de Montpellier II et Paris VII
\msp{page 116}
on trouve un foncteur
\[
i \longmapsto \underbrace{\pi_0(X_i)}_{\simeq\, \pi_0(X/i)} : I \longrightarrow \mathrm{Ens}.
\]
Ceci posé, on a le résultat sorital

\textbf{Proposition 12 :} \textit{Soit $f : X \longrightarrow I$ $W$-propre, et supposons que $W$ satisfasse l'axiome de connexité} W(5) (sinon, la $W$-propreté ne nous sera d'aucune utilité \dots) (\footnote[81]{Il vaut mieux dire : Soit $f$ $W_0$-propre, où $W_0$ est le localiseur de la 0-connexité \dots}). \textit{On a alors une bijection canonique}
\[
\pi_0(X) \xleftarrow{\ \sim\ } \varinjlim_{I} \pi_0(X_i), \tag*{$(*)$}
\]
\textit{où la flèche est induite par les applications canoniques $\pi_0(X_i) \longrightarrow \pi_0(X)$.}

\textsc{Démonstration.} Passant à $f^\circ : X^\circ \longrightarrow I^\circ$, on est ramené à prouver le même résultat pour $X$ \textit{$W$-lisse} sur $I$, avec la seule différence que le foncteur
\[
i \longmapsto \pi_0(X_i)
\]
sur $I$ sera contravariant, i.e. un préfaisceau, ce qui offre l'avantage de l'intuition attachée au formalisme des préfaisceaux. Soit
\[
g : I \longrightarrow e, \qquad \text{d'où } gf = p : X \longrightarrow e,
\]
il est immédiat, par le calcul standard
\[
p_!(F) = \pi_0(X/F) \qquad (F \in \mathrm{Ob}\, X^\wedge)
\]
(où on a identifié $e^\wedge$ à (Ens)), qu'on a
\[
\pi_0(X) = p_!(e_X) = g_!(f_!(e_X)).
\]
Or le préfaisceau
\[
i \longmapsto \pi_0(X_i) \qquad (\simeq \pi_0(i\backslash X), \text{ car } X\ W\text{-lisse sur } I)
\]
\msp{page 117}
n'est autre que $f_!(e_X)$. D'autre part, pour un préfaisceau quelconque $G$ sur $I$, on a
\[
g_!(G) \simeq \pi_0(I/G) \simeq \varinjlim_{I} G(i),
\]
et appliquant ceci à $f_!(e_X)$, on trouve la formule $(*)$ de la proposition 12.

Revenons à la situation avec $f$ \textit{$W$-propre,} et considérons un foncteur
\[
p : X \longrightarrow Y,
\]
nous sommes intéressés à ‘calculer’
\[
p_!(e_X) = (y \longmapsto \pi_0(y\backslash X)) \in \mathrm{Ob}\, Y^\wedge.
\]
Considérons pour ceci le diagramme suivant
% foot: Version 7 décembre 2002
\printedpage{80}
% head: A. Grothendieck : Dérivateurs, Ch. XII, Édition des Universités de Montpellier II et Paris VII
\[
\begin{tikzcd}
y\backslash X_i \arrow[r] \arrow[d] & y\backslash X \arrow[r] \arrow[d] & y\backslash Y \arrow[d] \\
X_i \arrow[r] \arrow[d] & X \arrow[r, "p"] \arrow[d, "f"] & Y \\
\underset{\simeq\, e}{\{i\}} \arrow[r] & I\rlap{\quad ,} &
\end{tikzcd}
\]
où les trois carrés sont cartésiens. Comme $y\backslash Y \longrightarrow Y$ est $W$-propre, il en est de même de $y\backslash X \longrightarrow X$ par changement de base, donc du composé $y\backslash X \longrightarrow X \longrightarrow I$ par transitivité. Le diagramme montre d'autre part que la fibre de $y\backslash X$ en $i \in \mathrm{Ob}\, I$ s'identifie à $y\backslash X_i$.

\msp{page 118}
La proposition 12 nous donne donc la formule
\[
\pi_0(y\backslash X) = \varinjlim_{I} \pi_0(y\backslash X_i),
\]
laquelle formule ne peut s'empêcher d'être fonctorielle en $y$, tout comme celle de prop. 12 l'est dans $\underline{\mathrm{Prop}}_W I$. Or pour $y$ variable dans $Y$, le préfaisceau sur $Y$
\[
y \longmapsto \pi_0(y\backslash X_i)
\]
n'est autre que $p_{i\,!}(e_{X_i})$, où
\[
p_i : X_i \longrightarrow Y
\]
est le composé $X_i \longrightarrow X \xrightarrow{p} Y$ (induit par $p$). On a ainsi prouvé

\textbf{Corollaire 1} (de prop. 12). \textit{Soit $f : X \longrightarrow I$ $W$-propre} (\footnote[82]{dire $W_0$-propre \dots}), \textit{et $p : X \longrightarrow Y$ un foncteur quelconque. Soit, pour tout $i \in \mathrm{Ob}\, I$, $p_i : X_i \longrightarrow Y$ le foncteur induit sur la fibre $X_i$ de $X$ en $i$. On a alors une loi} fonctorielle \textit{naturelle}
\[
i \longmapsto p_{i\,!}(e_{X_i}) : I \longrightarrow Y^\wedge,
\]
\textit{et on a la formule}
\[
p_!(e_X) \xleftarrow{\ \sim\ } \varinjlim_{I} p_{i\,!}(e_{X_i}), \tag*{$(*)$}
\]
\textit{où la flèche est induite par les morphismes canoniques}
\[
p_{i\,!}(e_{X_i}) \longrightarrow p_!(e_X)
\]
% foot: Version 7 décembre 2002
\printedpage{81}
% head: A. Grothendieck : Dérivateurs, Ch. XII, Édition des Universités de Montpellier II et Paris VII
\msp{page 119}
\textit{induites par les inclusions $X_i \longrightarrow X$.}

La loi fonctorielle de $p_{i\,!}(e_{X_i})$ en $i$ a été explicitée en cours de démonstration.

Notons que lorsque $X$ sur $I$ est cofibrée et scindée, i.e. provient d'un foncteur
\[
i \longmapsto X_i, \quad I \longrightarrow \mathrm{Cat},
\]
alors la loi fonctorielle de $p_{i\,!}(e_{X_i})$ en $i$ s'explicite de façon particulièrement agréable : à une flèche
\[
\alpha : i \longrightarrow j, \qquad \text{d'où } \alpha_* : X_i \longrightarrow X_j,
\]
est associé
\[
p_!(\alpha) : p_{i\,!}(e_{X_i}) \longrightarrow p_{j\,!}(e_{X_j}),
\]
en notant qu'on a un diagramme
\[
\begin{tikzcd}
X_i \arrow[rr, "\alpha_*"] \arrow[dr, "p_i"', ""{name=A}] & & X_j \arrow[dl, "p_j", ""'{name=B}] \\
& Y &
\arrow[from=A, to=B, bend left, dashed, "\lambda_\alpha"] % ?? uncertain: line style of the curved arrow lambda_alpha -- dashed or dotted
\end{tikzcd}
\]
avec une flèche de commutativité
\[
\lambda_\alpha : p_i \longrightarrow p_j\alpha_*
\]
(pas \texttt{[un]} isomorphisme en général), d'où
\[
p_{i\,!}(e_{X_i}) \xrightarrow{(\lambda_\alpha)_!} p_{j\,!}\underbrace{\alpha_{*\,!}(e_{X_i})}_{\longrightarrow\, e_{X_j}} \longrightarrow p_{j\,!}(e_{X_j}),
\]
plus de doute c'est lui! (\footnote[83]{Cette flèche \texttt{[la flèche }$(\lambda_\alpha)_!$\texttt{]} devrait être explicitée avec soin!}).

\msp{page 120}
\textbf{Corollaire 2 :} \textit{Sous les conditions du corollaire 1, supposons que $I$ ait un objet final $i_0$. Alors}
\[
(p_{i_0})_!(e_{X_{i_0}}) \longrightarrow p_!(e_X)
\]
\textit{est un isomorphisme dans $Y^\wedge$.}

D'où aussitôt la proposition 11 : si $X_{i_0}$ est totalement $W$-asphérique, alors $X$ l'est, car pour tout $p : X \longrightarrow Y$, $p_!(e_X) \simeq p_{i_0\,!}(X_{i_0})$, et ce dernier élément de $Y^\wedge$ est $W$-asphérique % [sic] p_{i_0 !}(X_{i_0}) -> p_{i_0 !}(e_{X_{i_0}}), elsewhere the argument of p_{i!} is e_{X_i}
par hypothèse sur $X_{i_0}$.

Notons le cas particulier :

\textbf{Corollaire 3.} \textit{Soit $Z$ un objet totalement $W$-asphérique dans} Cat. \textit{Alors son produit $Z\times I$ par toute $I$ dans} Cat \textit{ayant un objet final est aussi totalement $W$-asphérique.}
% foot: Version 7 décembre 2002
\printedpage{82}
% head: A. Grothendieck : Dérivateurs, Ch. XII, Édition des Universités de Montpellier II et Paris VII

\msp{page 121}

\section*{3\quad Sommes amalgamées et intégration\quad \begin{tikzcd}[ampersand replacement=\&, row sep=small, column sep=small, baseline=(\tikzcdmatrixname-2-1.base)]
\& X_1 \\
\int X_0 \arrow[ur] \arrow[dr] \& \\
\& X_2
\end{tikzcd}}

\subsection*{3.1\quad Réduction canonique au cas d'un diagramme d'immersions fermées}

Considérons un diagramme
\[
\begin{tikzcd}[row sep=small, column sep=small]
& X_1 \\
X_0 \arrow[ur, "\alpha_1"] \arrow[dr, "\alpha_2"'] & \\
& X_2
\end{tikzcd}
\tag*{$(*)$}
\]
dans Cat. Nous pouvons le regarder comme un contrafoncteur
\[
\Phi^\circ \longrightarrow \mathrm{Cat},
\]
où
\[
\Phi = \left( \vcenter{\hbox{\begin{tikzcd}[row sep=small, column sep=small]
& 1 \arrow[dl] \\
0 & \\
& 2 \arrow[ul]
\end{tikzcd}}} \right),
\]
donc comme définissant une catégorie fibrée (scindée) $\mathscr{X}$ sur $\Phi$. Il est parfois commode de représenter cette catégorie $\mathscr{X}$ par le ‘signe d'integration’ % [sic] integration -> intégration, accent missing in the print
\[
\begin{tikzcd}[row sep=small, column sep=small]
& X_1 \\
\int X_0 \arrow[ur] \arrow[dr] & \\
& X_2,
\end{tikzcd}
\]
et la regarder comme une sorte de ‘somme amalgamée homotopique’ de $X_1$ et $X_2$ sous $X_0$, ou comme une ‘limite inductive homotopique’ du diagramme de départ $(*)$.

\msp{page 122}
Il faut entendre par là que (pour $W \subseteq \mathrm{Fl}(\mathrm{Cat})$ donné) l'image $\mathrm{hot}_W(\mathscr{X})$ de $\mathscr{X}$ dans $\mathrm{Hot}_W$ a un ‘sens homotopique’ remarquable en termes de $(*)$, de façon précise, le diagramme $(*)$ définit un objet $\xi$ de $\mathrm{HOT}_W(\Phi)$ (qui est une catégorie de fractions de la catégorie de tous les diagrammes du type $(*)$ \dots), et si
\[
\pi : \Phi \longrightarrow e
\]
désigne la projection absolue de $\Phi$, on a
\[
\pi_!(\xi) = \mathrm{hot}_W(\mathscr{X}),
\]
% foot: Version 7 décembre 2002

\end{document}
